\chapter{Ideals and Dedekind Domains}\label{ch:ideals}

This chapter develops ideal theory in $\Od$,
following Boxer Lectures~8--11 and~14--17.

\textbf{Status: planned (long-term).}

% ----------------------------------------------------------------
\section{Dedekind Domain Structure}

\begin{theorem}[$\Od$ is a Dedekind domain]
  \label{thm:plan_dedekind}
  \notready
  \uses{thm:classification}
  For every squarefree $d\notin\{0,1\}$, the ring $\Od$ is a
  Dedekind domain: Noetherian, integrally closed, and every
  nonzero prime ideal is maximal.
\end{theorem}

\begin{remark}
  Mathlib already proves that $\OK$ is a Dedekind domain for any
  number field $K$ (\texttt{NumberField.RingOfIntegers.isDedekindDomain}).
  The formalization task is to transport this to the explicit
  models $\Z[\sqrt{d}]$ and $\Z[(1+\sqrt{d})/2]$ via the ring
  isomorphism.
\end{remark}

% ----------------------------------------------------------------
\section{Unique Factorization of Ideals}

\begin{theorem}[Unique factorization into prime ideals]
  \label{thm:plan_ideal_factorization}
  \notready
  \uses{thm:plan_dedekind}
  Every nonzero ideal $I\subseteq\Od$ factors uniquely as
  $I = \mathfrak{p}_1^{e_1}\cdots\mathfrak{p}_r^{e_r}$
  with $\mathfrak{p}_i$ prime ideals.
\end{theorem}

% ----------------------------------------------------------------
\section{Splitting of Rational Primes}

\begin{definition}[Splitting types]\label{def:plan_splitting}
  \notready
  A rational prime $p$ is:
  \begin{itemize}
    \item \emph{split} if $p\Od = \mathfrak{p}\bar{\mathfrak{p}}$
      with $\mathfrak{p}\ne\bar{\mathfrak{p}}$;
    \item \emph{inert} if $p\Od$ is prime;
    \item \emph{ramified} if $p\Od = \mathfrak{p}^2$.
  \end{itemize}
\end{definition}

\begin{theorem}[Splitting criterion for quadratic fields]
  \label{thm:plan_splitting_criterion}
  \notready
  \uses{thm:plan_disc_quadratic, def:plan_splitting}
  Let $p$ be a rational prime, $d$ squarefree.
  \begin{enumerate}
    \item $p$ ramifies iff $p\mid\disc(\Od)$.
    \item For $p\nmid\disc(\Od)$: $p$ splits iff
      $\bigl(\frac{d}{p}\bigr) = 1$, and $p$ is inert iff
      $\bigl(\frac{d}{p}\bigr) = -1$.
  \end{enumerate}
\end{theorem}

% ----------------------------------------------------------------
\section{Class Number}

\begin{definition}[Ideal class group]\label{def:plan_class_group}
  \notready
  \uses{thm:plan_ideal_factorization}
  The \emph{class group} $\mathrm{Cl}(\Od)$ is the quotient of
  fractional ideals by principal ideals.  The \emph{class number}
  is $h_d := |\mathrm{Cl}(\Od)|$.
\end{definition}

\begin{remark}
  Computing $h_d$ for specific $d$ (e.g.\ showing $h_{-5}=2$,
  $h_{-163}=1$) would be a natural verification target.
  Mathlib provides \texttt{ClassGroup} infrastructure.
\end{remark}

% ----------------------------------------------------------------
\section{Formalization Notes}

\begin{remark}
  This chapter is a longer-term goal.  The recommended path is:
  \begin{enumerate}
    \item Transport \texttt{IsDedekindDomain} from
      $\OK$ to $\Z[\sqrt{d}]$ / $\Z[(1+\sqrt{d})/2]$.
    \item Use Mathlib's
      \texttt{IsDedekindDomain.exists\_prime\_ideal\_of\_dvd}
      to exhibit factorizations.
    \item For splitting, connect with
      \texttt{NumberField.RingOfIntegers.splits\_iff}.
  \end{enumerate}
\end{remark}
